\section{Appendix: Scope and Qualifications}

This study treats the 1962 \work{Missale Romanum}'s Twentieth Sunday after
Pentecost, second class, green, printed pp.~412--413, for the requested
general-calendar occurrence on 11 October 2026. The FSSP Province of
France's dated Ordo corroborates that Sunday within its institutional
scope. No local patronal or dedication calendar was supplied. The study
does not determine current permission to celebrate in a particular church.
The Maternity commemoration's three orations are included conditionally
at Low or conventual Mass and omitted at a nonconventual sung Mass. The
Trinity Preface remains in both branches. The Ordinary beyond that
appointed Preface is not reproduced.

The normalized Latin was collated in the research stage against the 1962
facsimile and the public-domain antecedents in Pustet's 1862 Missal and
Cummiskey's 1861 Preface. The facsimile itself is not treated as a
blanket public-domain book. Cummiskey supplies the historical Sunday
prayer English; Lasance--Walsh supplies the three Maternity translations
under the recorded United States nonrenewal basis, outside the excluded
later inserts. Their English is historical study material, not an
approved 1962 English liturgical edition. Biblical English comes from
the Douay--Rheims--Challoner Gutenberg witness; its ebook number is an
identifier, not a printing year. Adapted Latin chants remain distinct
from their biblical English context. No new prayer translation is offered.

The reception sweep is bounded by the checked works and loci below.
Augustine, Gregory, Chrysostom, Jerome, and Bellarmine comment on the
constituent passages; the two complete-formulary readings are editorial
syntheses. Schuster expounds this Sunday's matching base elements as a
Mass. He does not expound this dated Marian commemoration or the Trinity
Preface in that chapter. No direct patristic commentary on those composed
forms was established. Patristic and saintly spiritual interpretation
does not date a text's composition. Historical translations and modern
transcriptions are not newly established critical editions.

The historical appendix preserves received attribution, composition
alternatives, derived event dating, and explicit gaps. Its Psalter
boundary is not a separate composition estimate for every psalm.
Chronological reference works and modern Catholic biblical introductions
remain scholarly witnesses, not doctrinal definitions of a year.
The exact text and evidence records are \texttt{propers/verified.md},
\texttt{research/scope.md}, and \texttt{research/interpretations.md}.
This is an AI-assisted Catholic study aid, not an official liturgical
book or an ecclesiastical judgment.

\section{References}
\begin{description}[style=nextline,leftmargin=0pt,labelsep=0pt,itemsep=.45em]
\item[Missal and occurrence]
\work{Missale Romanum}, Vatican typical edition, 1962, pp.~412--413,
nn.~1689--1698; pp.~685--686, nn.~3847, 3855, 3857; Trinity Preface,
pp.~244--245, n.~1053. General rubrics, nn.~91,106--109,111, and
Mass rubrics, nn.~427--429,433--437,480,483,494,505.
CMAA facsimile: \url{https://media.churchmusicassociation.org/pdf/missale62.pdf}.
FSSP Province de France, \work{Ordo du mois}, October 2026, 11 October
row, consulted 5 October 2026: \url{https://www.fssp.fr/ordo-du-mois/}.

\item[Historical liturgical witnesses]
\work{Missale Romanum}, Pustet, Ratisbon, 1862, pp.~352--354 and
appendix [171]--[172]. \work{Roman Missal Translated into the English
Language for the Use of the Laity}, first revised edition, Eugene
Cummiskey, Philadelphia, 1861, ``Twentieth Sunday after Pentecost,''
pp.~448--450; Prefaces, xxviii--xxix:
\url{https://archive.org/details/romanmissaltran00churgoog}.
Francis Xavier Lasance and Francis Augustine Walsh, \work{The New Roman
Missal}, revised 1945 edition, later printing, ``Maternity of the Blessed
Virgin Mary,'' pp.~1233, 1235--1236. The historical wording retains its
own public-domain status in the United States; it is not project-owned prose.

\item[Scripture]
Douay--Rheims--Challoner, Project Gutenberg eBook 1581: Daniel~3;
Psalms~107,118,136,144; John~4; Ephesians~4--6; Jeremiah~25.
\url{https://www.gutenberg.org/ebooks/1581}.
Psalm references use Vulgate numbering, with common modern numbers supplied
where useful. Chant adaptations and appointed verse cuts are identified
beside their study texts.

\item[Augustine]
\work{Tractates on John}, XVI.1--7, NPNF I.7; \work{Expositions of the
Psalms}, Ps.~107.1--2; Ps.~118, \S\S1--8,50--56; Ps.~136.1--13;
Ps.~144.14--18, NPNF I.8. The English psalm headings are 108,119,137,145.
Historical English in the retained CCEL witnesses:
\url{https://www.ccel.org/ccel/schaff/npnf107.html};
\url{https://www.ccel.org/ccel/schaff/npnf108.html}.

\item[Gregory the Great]
\work{Homilies on the Gospels}, II.28.1--3, Latin homily for Nereus
and Achilleus, on John~4:46--53. Wikisource transcription,
\url{https://la.wikisource.org/wiki/Homiliarum_in_Evangelia}.
Modern transcription credited to Wikisource contributors, including
Mizardellorsa, under CC BY-SA 3.0. The exposition here paraphrases
Gregory's argument; no modern English translation is quoted.

\item[Chrysostom and Jerome]
John Chrysostom, \work{Homilies on Ephesians}, XIX, on 5:15--21,
including the final exhortation to reciprocal service, NPNF I.13;
the separate editorial note~404 reports Meyer's restriction of ``all
things'' to blessings. \url{https://www.ccel.org/ccel/schaff/npnf113.html}.
Jerome, \work{Commentary on Daniel}, I, at 3:23--29, especially 3:29,
PL~25, cols.~639--640, historical Latin transcription.

\item[Later saintly and liturgical reception]
Robert Bellarmine, \work{A Commentary on the Book of Psalms}, John
O'Sullivan's historical translation (1866), Ps.~136:1--4 and
144:14--17; \url{https://www.ecatholic2000.com/bellarmine/commentary-on-psalms.shtml}.
Ildefonso Schuster, \work{The Sacramentary}, translated by Arthur
Levelis-Marke, vol.~III, London: Burns Oates \& Washbourne, 1927,
``Twentieth Sunday after Pentecost / Quinta post natale Sancti Cypriani,''
pp.~174--178.

\item[Scriptural chronology and historical attribution]
George Leo Haydock, \work{Douay--Rheims Bible with Commentary},
Loreto Feeney memorial edition, 2014, historical notes at Ps.~118:1
(p.~781), Ps.~136 (p.~794), and Ps.~70 (p.~740, reporting Ussher).
\work{Catholic Encyclopedia}: Corbett, ``David, King,'' opening
chronology; Gigot, ``Book of Daniel,'' concluding section
``Deutero-canonical portions''; Ladeuze, ``Epistle to the Ephesians,''
``Date and place of composition''; Prat, ``St. Paul,'' chronological
summary; Fonck, ``Gospel of St. John,'' date discussion;
Durand, ``New Testament,'' Johannine date paragraph;
Maas, ``Chronology of the Life of Jesus Christ,'' ``Second journey''
and chronological calibration. The traditional dates are historical
scholarly proposals, with alternatives shown in the appendix.

\item[Jerusalem background and modern comparisons]
\work{Catholic Encyclopedia}: Reid, ``Captivities of the
Israelites,'' destruction paragraph; Meistermann, ``Temple of
Jerusalem,'' opening; Schets, ``Third and Fourth Books of Kings,''
authorship discussion; Sloet, ``Chronology of the Kings,'' Petavius
table, retained transcription (printed page unverified).
USCCB, NABRE introductions to Psalms and Ephesians, consulted 5 October
2026, especially the Psalter's dating qualification and Ephesians'
authorship/destination paragraphs:
\url{https://bible.usccb.org/bible/psalms/0};
\url{https://bible.usccb.org/bible/ephesians/0}.
These contemporary introductions are summarized, not reproduced; their
copyright does not pass to the project's prose license.
\end{description}
